\par \Proof Так как $\alpha$ алгебраичен над $F$, $m_{\alpha, F}=\sum_{i=0}^n a_i \alpha^i=0$. Тогда $\varphi(m_{\alpha, F})=\sum_{i=0}^n \psi(a_i) \psi(\alpha)^i=\sum_{i=0}^n \varphi(a_i) \cdot \beta^i = 0$. Существование очевидно.
\par Покажем единственность: предположим, что $\varphi$ тоже удовлетворяет требованиям в условии, тогда получим, что $\sum_{i=0}^n \varphi(a_i) \cdot \beta^i = 0 \Rightarrow$ что получили два многочлена минимальной степени, порождающих идеал вида $I=\{f(x) \in F'(\beta)[x] \: | \: f(\beta)=0\}$. Заметим, что этот идеал простой, а значит $(\psi(m_{\alpha, F})) \sim (\varphi(m_{\alpha, F}))$, но это значит, что $\varphi, \psi$ переводят один и тот же многочлен в тот же самый, а значит и все многочлены в те же самые, а значит они совпадают. Противоречие, из которого вытекает
единственность. \EndProof